\section{One Way Door: señales de un producto acertado a largo plazo}

Tras la discusión sobre Miaoya (妙鸭) y el diseño del contexto, hace falta un criterio para seleccionar productos. Zhang Yueguang (张月光) usa One Way Door para juzgar si vale la pena apostar a largo plazo por algo.

A Zhang Yueguang le gustan los productos One Way Door: una vez que el usuario adopta la nueva solución, sabe que ya no puede volver a la anterior. Cursor y Claude Code, frente a los IDE tradicionales, son un ejemplo típico. Un One Way Door no necesariamente ofrece una experiencia perfecta desde el principio, pero deja ver que en el futuro formará una ventaja integral.

\begin{figure}[H]
\centering
\includegraphics[width=0.92\textwidth]{one-way-door-product.png}
\caption{Criterio de producto One Way Door: solución anterior, nueva solución, nueva normalidad.}
\end{figure}

\begin{knowledgebox}{Lectura de la figura: la puerta de un solo sentido no es un wow moment}
En la figura, la puerta de un solo sentido se ubica entre la solución anterior y la nueva normalidad. Un wow moment solo deslumbra, y el usuario puede volver pronto al flujo anterior; la puerta de un solo sentido, en cambio, modifica el comportamiento predeterminado del usuario. Su participación de mercado puede empezar siendo pequeña, pero crece de forma sostenida con el tiempo.
\end{knowledgebox}

\subsection{Por qué un Agent puede ser una puerta de un solo sentido}

Si un producto Agent realmente permite al usuario superar los límites de su capacidad, no solo aumenta la eficiencia, sino que cambia la idea que el usuario tiene de «qué puedo hacer». Zhang Yueguang explica que Dokie entra por ahora a través de las PPT, pero que su verdadero objetivo no son las PPT, sino un amigo de IA que le ayude a superar los límites de su capacidad personal.

\begin{importantbox}{Prueba de producto de la puerta de un solo sentido}
Para saber si vale la pena desarrollar un producto de IA a largo plazo, pueden plantearse estas preguntas: después de usarlo, ¿no quiero volver al método anterior? ¿Resuelve una necesidad existente, pero de una manera nueva que ofrece una ventaja integral? ¿Se volverá cada vez más potente conforme avancen los modelos?
\end{importantbox}

\subsection{Resumen de la sección}

One Way Door es el filtro de productos de este episodio. Miaoya fue un gran éxito viral, pero no necesariamente una puerta de un solo sentido; un Agent, si logra superar de manera estable los límites de capacidad, puede convertirse en una puerta acertada a largo plazo.
